%% ncolumn.tex -- documentation for the ncolumn package
%%
%% Copyright (c) 2026 Munehiro Yamamoto <munepixyz@gmail.com>
%% License: LPPL 1.3c or later
%%
%% Build: pdflatex ncolumn.tex (twice)
\documentclass{article}
\usepackage[margin=28mm]{geometry}
\usepackage{xcolor}
\usepackage{ncolumn}
\usepackage[colorlinks, allcolors=blue!50!black]{hyperref}
\newcommand*\ncolversion{v0.1 (2026/11/23)}
\newcommand*\cs[1]{\texttt{\char`\\#1}}
\newcommand*\pkg[1]{\textsf{#1}}
\newcommand*\key[1]{\texttt{#1}}
\newcommand*\meta[1]{$\langle$\textit{#1}$\rangle$}
\newcommand*\marg[1]{\texttt{\char`\{}#1\texttt{\char`\}}}
\emergencystretch=2em
%% \filler{<n>}: <n> filler sentences
\makeatletter
\newcommand\filler[1]{%
  \count@=0
  \loop\advance\count@ by 1
    This sentence only fills the column, so that the page is long
    enough to show where the material goes.
  \ifnum\count@<#1 \repeat}
\makeatother
%% \demo{<height>}{<text>}: a grey box of <height> as a float body
\newcommand\demo[2]{\centering
  \fcolorbox{black!40}{black!8}{\parbox[c][#1][c]%
    {\dimexpr\linewidth-2\fboxsep-2\fboxrule}{\centering\small #2}}}

\title{The \pkg{ncolumn} package\\
  \large N columns in the way of \cs{twocolumn}, with material that
  spans columns}
\author{Munehiro Yamamoto}
\date{\ncolversion}

\begin{document}
\maketitle

\section{Introduction}

\pkg{multicol} collects the material of all columns and splits it into
the columns afterwards.  Ordinary floats cannot be used in it, and
packages that look at the room that is left on the page, such as
\pkg{tcolorbox} with \key{breakable}, need extra care.

\cs{twocolumn} of \LaTeX{} works in another way: it builds the columns
one after another.  A column is a page for \TeX, and the finished
columns are put side by side when the last one is done.  So
\texttt{figure} and \texttt{table}, footnotes, \cs{marginpar} and such
packages work in it as on a page.

The \pkg{ncolumn} package makes the output routine of \cs{twocolumn}
work for three, four or more columns, and adds two things:
\begin{itemize}
\item \emph{Spanning material.}  Floats, headings and boxes can span
  $k$ columns from one of the four corners of the page.
\item \emph{Footnotes.}  The footnotes of a page are put into each
  column, into the last column, or into an area of $M$ columns across
  the page.
\end{itemize}
With two columns and without the new keys, the pages are those of
\cs{twocolumn}.  Section~\ref{sec:examples} shows two pages that are
typeset with the package.

\paragraph{Supported formats/engines}
pdf\LaTeX, Xe\LaTeX, Lua\LaTeX{} (with or without Lua\TeX-ja),
p\LaTeX{} and up\LaTeX{} (with dvipdfmx); horizontal (yoko) and
vertical (tate) writing.  The package requires \LaTeX{} 2025-06-01 or
later.

\section{Usage}

\begin{verbatim}
\documentclass{article}
\usepackage[columns=3]{ncolumn}
\end{verbatim}
\key{columns=}\meta{N} sets the whole document in \meta{N} columns.
If the package is loaded without \key{columns} in a document whose
class has the option \key{twocolumn}, the document has two columns as
before.

\subsection{Changing the number of columns}

\begin{verbatim}
\ncolumn{4}
\ncolumn{3}[\section{A heading across all columns}]
\onecolumn
\twocolumn
\end{verbatim}
\cs{ncolumn}\marg{\meta{N}} starts a new page with
\meta{N} columns.  The material in its optional argument is put across
all columns at the top of this page, as with \cs{twocolumn}.
\cs{twocolumn} does not mean ``two columns'' but ``the columns of this
document'': it goes back to the number of columns that was used
before.  The classes call it after \cs{maketitle} or an index to go
back to the columns.

\subsection{Spanning floats}

\begin{verbatim}
\begin{figure*}[t]               % all columns, at the top
\begin{figure*}[b]               % all columns, at the bottom
\begin{figure*}[span=2,at=tl]    % two columns from the first one, top
\begin{figure*}[span=2,at=br]    % two columns up to the last one, bottom
\begin{figure*}[!,span=2,at=tr]  % ! as for ordinary floats
\end{verbatim}
An argument of \texttt{figure*} or \texttt{table*} with a \texttt{=}
is read as a list of keys.  \key{span} is the number of columns and
\key{at} is the corner: \key{tl}, \key{tr}, \key{bl} or \key{br}.
\texttt{[b]}, which \LaTeX{} does not allow for these floats, puts a
float across all columns at the bottom.

The corner is given in the directions of the text: \key{t} and \key{b}
are the top and the bottom of a column, \key{l} and \key{r} are the
side of the first and of the last column.  In yoko this is what it
looks like.  In tate the first column is the upper one and the top of
a column is on the right, so \key{tl} is the upper right corner of
the page.

A page can have several pieces of spanning material.  A piece is put
on a page only if the columns it covers have the same height reserved
so far; otherwise it would leave a hole, and it goes to the next page.

As in \LaTeX, a spanning float is put on a later page than the one
where it is met.  With the key \key{same-page-spans}, it is put on the
page where it is met
\begin{itemize}
\item at the top, if none of the columns it covers is begun;
\item at the bottom, if it covers only the current column and columns
  after it, and the current column has room.
\end{itemize}
No spanning float is put at the top of a page that begins with the
material of \cs{ncolumn}\marg{\meta{N}}\texttt{[...]} or with a
\cs{spanblock} at the top, as with \cs{twocolumn}\texttt{[...]} in
\LaTeX.

In one column, \key{span} and \key{at} are ignored and the float is an
ordinary one.

\subsection{Spanning blocks}

\begin{verbatim}
\spanblock[span=3,at=tl]{\section{A large heading}}
\spanblock*[span=2,at=br]{...}
\end{verbatim}
\cs{spanblock} puts material that is not a float, such as a heading or
a \pkg{tcolorbox}, across columns.  Without the star it starts a new
page and puts the block on it; with the star the block is put on a
later page, like a float.

\subsection{Footnotes}

\begin{verbatim}
\usepackage[columns=3,footnotes=last-column]{ncolumn}
\ncolumnsetup{footnotes=span,footnote-columns=2}
\end{verbatim}
The key \key{footnotes} has one of these values:
\begin{description}
\item[\key{column}] The footnotes are put at the bottom of the column
  of their marks, as in \LaTeX{} (default).
\item[\key{last-column}] All footnotes of a page are put at the bottom
  of its last column, as \pkg{ftnright} does for two columns.  One run
  is enough.
\item[\key{span}] The footnotes are put into an area across all
  columns at the bottom of the page.  With
  \key{footnote-columns=}\meta{M}, this area has \meta{M} columns.
  Several runs are needed.
\end{description}
With \key{span}, the columns that are finished cannot be made shorter
for a footnote of a later column.  So the height of the footnote area
of each page is written to the \texttt{.aux} file and is taken off all
columns in the next run.  The package warns to run again when a height
has changed.

Reserving the room can move the line of the mark to the next page, and
not reserving it brings it back.  If the package finds this, it keeps
the room reserved (the page is then shorter by a footnote area that it
does not have).  Up to about five runs can be needed.
\key{footnote-reserve=}\meta{length} reserves a fixed height on every
page and does not use the \texttt{.aux} file.

In tate with \pkg{jlreq}, \key{last-column} is the default, as
\pkg{jlreq} puts the footnotes at the bottom of the last column.

\subsection{Marginal notes}

With three or more columns, a \cs{marginpar} goes to the nearer
margin: to the margin on the side of the first column from the first
half of the columns, to the other margin from the second half.  The
notes of columns that share a margin do not overlap.

\subsection{\pkg{tcolorbox}}

A \pkg{tcolorbox} with \key{breakable} is broken across columns as it
is across pages.  Two packages for \pkg{tcolorbox} are tested with
\pkg{ncolumn}:
\begin{itemize}
\item \pkg{tcbpagefootnote} (v1.4 or later): with the key \key{page
  footnote}, the footnotes of a box are footnotes of the text.
  \key{footnotes} puts them as it puts the others, for each part of a
  box that is broken across columns, and in a tagged PDF too.
\item \pkg{tcbnestbreak} (v0.2): breakable boxes inside a breakable
  box are broken across columns as well.
\end{itemize}
A box in spanning material cannot have footnotes of the text (see
section~\ref{sec:limitations}).

\subsection{Options and parameters}

\begin{center}
\begin{tabular}{@{}lp{95mm}@{}}
\hline
\key{columns=}\meta{N} & the number of columns of the document
  (\cs{usepackage} only) \\
\key{footnotes=}\meta{value} & where the footnotes are put:
  \key{column}, \key{last-column} or \key{span} \\
\key{footnote-columns=}\meta{M} & the columns of the footnote area
  with \key{span} \\
\key{footnote-reserve=}\meta{length} & a fixed height of the footnote
  area with \key{span} \\
\key{same-page-spans} & put spanning floats on the page where they are
  met, if possible \\
\hline
\end{tabular}
\end{center}
The keys are given to \cs{usepackage} or, but \key{columns}, to
\cs{ncolumnsetup}.

The parameters of \LaTeX{} for two columns are used as they are:
\cs{columnsep}, \cs{columnseprule}, \texttt{dbltopnumber},
\cs{dbltopfraction}, \cs{dblfloatsep}, \cs{dbltextfloatsep} and
\cs{dblfigrule}.  For spanning material at the bottom, there are the
counter \texttt{dblbotnumber} (default 2) and \cs{dblbotfraction}
(default 0.5).

\section{Examples}\label{sec:examples}

The next two pages are typeset with \key{same-page-spans} and with a
rule between the columns (\cs{columnseprule}).
\begin{itemize}
\item The first page has three columns.  Figure~\ref{fig:col} is an
  ordinary \texttt{figure} with \texttt{[t]}.  Figure~\ref{fig:tr} is
  a \texttt{figure*} with \texttt{[span=2,at=tr]} and
  Figure~\ref{fig:bl} one with \texttt{[span=2,at=bl]}; both are given
  at the beginning of the first column.  The footnotes are at the
  bottom of their columns (\key{footnotes=column}).
\item The second page has four columns.  The text with the frame is a
  \cs{spanblock*} with \texttt{[span=2,at=tr]} and Figure~\ref{fig:b}
  is a \texttt{figure*} with \texttt{[b]}; both are given at the
  beginning of the first column.  All footnotes of the page are in its
  last column (\key{footnotes=last-column}).
\end{itemize}

%% ------------------------------------------------------------------
\ncolumn{3}
\ncolumnsetup{same-page-spans}
\setlength\columnseprule{.4pt}
\begin{figure}[t]
\demo{12mm}{\texttt{figure} with \texttt{[t]}}
\caption{A float of one column}\label{fig:col}
\end{figure}
\begin{figure*}[span=2,at=tr]
\demo{20mm}{\texttt{[span=2,at=tr]}: two columns up to the last one,
  at the top}
\caption{Two columns from the top right corner}\label{fig:tr}
\end{figure*}
\begin{figure*}[span=2,at=bl]
\demo{20mm}{\texttt{[span=2,at=bl]}: two columns from the first one,
  at the bottom}
\caption{Two columns from the bottom left corner}\label{fig:bl}
\end{figure*}
\filler{4}\footnote{A footnote of the first column.}
\filler{22}\footnote{A footnote of a later column.}
\filler{10}

\ncolumn{4}
\ncolumnsetup{footnotes=last-column}
\spanblock*[span=2,at=tr]{\fbox{\parbox{\dimexpr\linewidth-2\fboxsep-2\fboxrule}{%
  \small\textbf{A block across two columns.}  \cs{spanblock} takes
  material that is not a float: a heading, a framed text like this
  one, or a \pkg{tcolorbox}.}}}
\begin{figure*}[b]
\demo{14mm}{\texttt{[b]}: all columns, at the bottom}
\caption{All columns at the bottom}\label{fig:b}
\end{figure*}
\filler{3}\footnote{A footnote from the first column.}
\filler{12}\footnote{A footnote from a later column.}
\filler{24}

\onecolumn
\setlength\columnseprule{0pt}

\section{How it works}

\begin{enumerate}
\item \emph{N columns.}  \cs{@outputdblcol} keeps the finished columns
  in $N$ boxes instead of one.  \cs{if@twocolumn} stays true and
  \cs{if@firstcolumn} is true in the first column only, so code that
  tests them works as before.
\item \emph{Spanning material.}  Each column has a height that is
  reserved at its top and one at its bottom, and the height of the
  column (\cs{@colht}) is set for each column.  The spanning material
  is laid over the columns when the page is put together.
\item \emph{Footnotes.}  \cs{@makecol} and \cs{@outputpage} are not
  redefined: p\LaTeX{} has its own versions for tate, Lua\TeX-ja uses
  their hooks, and so does the tagging code of \LaTeX{}
  (\cs{DocumentMetadata}).  The package takes \verb|\box\footins| at
  the beginning of \cs{output} instead.
\item \emph{Marks.}  The code of \texttt{ltmarks} that updates the
  marks for two columns is replaced by one for $N$ columns (the marks
  of each column are also in the regions \texttt{mcol-1},
  \texttt{mcol-2}, \ldots).  This is the only place where internal
  variables of another module are used.
\end{enumerate}
The internal macros of \LaTeX{} that are redefined or wrapped are
listed at the top of \texttt{ncolumn.sty}.

\section{Limitations}\label{sec:limitations}

\begin{itemize}
\item Material spans columns from the corners of the page only.  It
  cannot be put in the middle of a page with text above and below it:
  the columns would have to be balanced for that.  The columns of the
  last page are not balanced either.
\item When the floats are put out by \cs{clearpage}, a float that
  spans some of the columns is not put into its corner but on a float
  page, at the side of the first or of the last column.
\item With \key{footnotes=span}, a footnote is not split across pages.
\item A footnote inside spanning material is not typeset, as in an
  ordinary float or in \verb|\twocolumn[...]|.  Put \cs{footnotemark}
  into the material and \cs{footnotetext} after it.
\item In tate with \pkg{jlreq}, \key{footnotes=column} and
  \key{footnotes=span} do not work with the code of \pkg{jlreq} for
  its notes, which is written for two columns.  Use \key{last-column}.
\item These packages redefine the same macros; loading one of them
  with \pkg{ncolumn} is an error: \pkg{ftnright}, \pkg{nidanfloat},
  \pkg{dblfloatfix}, \pkg{cuted}, \pkg{balance}, \pkg{flushend},
  \pkg{pbalance}, \pkg{flowfram}, \pkg{dblfnote}.  \pkg{stfloats}
  gives a warning only, as \pkg{jlreq} loads it with Lua\LaTeX;
  \pkg{ncolumn} replaces its placement of the spanning floats.
\end{itemize}

\section{Tests}

\texttt{make test} compiles the test files \texttt{test-*.tex} with
pdf\LaTeX, Xe\LaTeX, Lua\LaTeX, p\LaTeX{} and up\LaTeX.  For each page
that is shipped out, \texttt{test-common.tex} records the state of the
package: the heights that are reserved at the top and at the bottom of
each column, the spanning material, and the first and the last mark of
each column.  At \verb|\end{document}| it compares this with the state
that the test file expects, reads the \texttt{.aux} file back and
checks that every footnote text is on the page of its mark, and writes
\texttt{NCOL-RESULT OK} or \texttt{NCOL-RESULT NG} to the log.

\texttt{test-twocolumn.tex} is compiled with and without the package:
with two columns, every page is the same, pixel by pixel, as with
\cs{twocolumn} of the \LaTeX{} kernel.
\texttt{test-tcbpagefootnote.tex} and \texttt{test-tcbnestbreak.tex}
also check that every footnote text has the width of the footnotes of
the page and that no footnote is split or held over.
\texttt{test-text.tex},
\texttt{test-japanese.tex} (yoko and tate) and
\texttt{test-tcolorbox.tex} are documents with real text that are also
meant to be looked at.  See \texttt{README.md} for the list of the
tests.

\section{History}

\begin{description}
\item[v0.1 (2026/11/23)] First version.
\end{description}

\section{License}

Copyright \copyright\ 2026 Munehiro Yamamoto.

This package is licensed under the conditions of the \LaTeX{} Project
Public License, either version 1.3c of this license or (at your
option) any later version.  The latest version of this license is in
\url{https://www.latex-project.org/lppl.txt}.

\end{document}
